\section{Feedforward Neural Network}
\label{sec:fnn}

% TODO: fill this section once FNN experiments are complete (see 3_dl.ipynb / 2_fnn.ipynb)

\subsection{Architecture}

% TODO: describe the FNN architecture (number of layers, hidden units, activation
%       functions, dropout, weight initialization scheme used).
%
% Key points to cover:
%   - Input layer: 10 lag features
%   - Hidden layers: [describe configuration]
%   - Activation: ReLU (automatically solves symmetry blindness)
%   - Dropout: [rate]
%   - Weight init: He initialization
%   - Output: single neuron (no activation), predicts absolute return

\subsection{Training Setup}

% TODO: describe training procedure.
%
%   - Loss function: [MSE / MAE / custom?]
%   - Optimizer: [SGD / Adam?], learning rate schedule
%   - Batch size, number of epochs
%   - Early stopping patience
%   - Train / validation / test split

\subsection{Results}

% TODO: insert results table (MAE, RMSE, QLIKE — full, calm, crisis).
%
% Preliminary reference values from notebook:
%   Custom NumPy FNN  — MAE: 0.605, val loss: 0.759
%   PyTorch reference — MAE: 0.579

% TODO: figure placeholder — predicted vs actual (FNN)

% TODO: figure placeholder — training / validation loss curves

\subsection{Analysis}

% TODO: discuss weight visualizations, gradient norms, lag sensitivity.
%       How does varying the number of lags affect performance?
%       SHAP / LIME feature importance (if computed).

\subsection{Limitations}

% TODO: describe remaining limitations of FNN:
%   1. Order blindness — input is a flat vector; temporal ordering is ignored.
%   2. Regime blindness — still a single model with no state.
%
% These motivate the RNN architecture in Section~\ref{sec:rnn}.
